\chapter{OBJETIVOS}

% ===========================================================================
% Este capítulo recoge el objetivo general, los objetivos específicos, la
% hipótesis de trabajo, el alcance y restricciones, la metodología seguida y
% la planificación temporal. Sustituye los textos de ejemplo por los tuyos.
% (Cubre los ítems IT-44 hipótesis/restricciones, IT-45 requisitos, IT-48 Gantt)
% ===========================================================================

\section{Objetivo general}

Diseñar, implementar y evaluar un sistema conversacional capaz de responder, a
partir de información oficial y actualizada, preguntas de estudiantes
preuniversitarios sobre las titulaciones de la Escuela Politécnica Superior de
Jaén, empleando técnicas de \textit{Retrieval-Augmented Generation} (RAG) sobre
un modelo de lenguaje de código abierto.

\section{Objetivos específicos}

% Deben ser verificables: al final del TFG se ha de poder decir sí/no para cada uno.
\begin{itemize}
    \item OE1. Construir un \textit{corpus} estructurado y trazable de las guías
    docentes, planes de estudio y salidas profesionales de la EPSJ mediante
    \textit{web scraping}.
    \item OE2. Definir y evaluar experimentalmente una estrategia de
    segmentación (\textit{chunking}) del corpus adecuada al dominio.
    \item OE3. Seleccionar, con criterios cuantitativos, el modelo de
    \textit{embeddings} y la base de datos vectorial del sistema.
    \item OE4. Implementar un \textit{pipeline} RAG completo sobre un LLM de
    código abierto desplegado en local.
    \item OE5. Evaluar el sistema con métricas objetivas de recuperación y de
    generación, y validarlo con usuarios reales.
    \item OE6. Ofrecer el sistema a través de una aplicación web de chat
    accesible sin registro.
\end{itemize}

\section{Hipótesis de trabajo}

% IT-44. Debe ser falsable.
Es posible construir, con recursos de cómputo modestos y modelos de código
abierto, un sistema RAG que responda preguntas sobre la oferta académica de la
EPSJ con una fidelidad (\textit{faithfulness}) y una capacidad de recuperación
(\textit{Recall@K}) suficientes para resultar útil a un estudiante
preuniversitario, sin incurrir en respuestas inventadas sobre asignaturas o
titulaciones.

\section{Alcance y restricciones}

% IT-44 / IT-45.
El sistema se limita a las titulaciones de grado de la EPSJ; no abarca el
conjunto de la Universidad de Jaén ni los estudios de posgrado. El proyecto se
desarrolla de forma individual, con una dedicación equivalente a los créditos
ECTS del TFG, y dentro del calendario de la convocatoria Extraordinaria~II del
curso 2025/26, con entrega el 3 de septiembre de 2026 y defensa el 9 de
septiembre de 2026.

Como limitación conocida del dominio, algunas asignaturas de cursos superiores
de titulaciones de implantación reciente no disponen todavía de guía docente
publicada, lo que se documenta explícitamente en los resultados.

\section{Metodología de trabajo}

% Justifica Kanban frente a Scrum para un TFG individual.
El desarrollo sigue una metodología ágil basada en \textbf{Kanban} con límite de
trabajo en curso (WIP), gestionada en \textit{GitHub Projects}. Se descarta
Scrum porque sus ceremonias y \textit{sprints} de duración fija están pensados
para equipos, mientras que este TFG lo desarrolla una sola persona; Kanban
permite un flujo continuo y un límite de WIP que evita la dispersión.

Cada ítem del \textit{backlog} se corresponde con una incidencia (\textit{issue})
en GitHub y, cuando implica una decisión de diseño relevante, con un registro de
decisión de arquitectura (ADR) versionado junto al código. Un ítem no se
considera terminado hasta que cumple su Definición de Hecho: código, prueba
asociada, ADR si procede y la sección de memoria correspondiente redactada. Esta
trazabilidad completa (requisito, \textit{issue}, código, prueba y sección de
memoria) es un principio rector del proyecto.

\section{Planificación temporal}

% IT-48. El Gantt refleja las fases reales y la fecha de entrega.
El trabajo se organiza en seis fases, sincronizadas con los \textit{milestones}
del repositorio. La Figura~\ref{fig.gantt} muestra su distribución temporal
hasta la entrega.

\begin{figure}[htbp]
    \centering
    \resizebox{\textwidth}{!}{%
    \begin{ganttchart}[
        hgrid, vgrid,
        bar/.append style={fill=blue!40},
        bar height=0.5,
        title height=1,
        y unit chart=0.7cm,
        x unit=0.115cm
        ]{1}{72}
        \gantttitle{Junio}{7}
        \gantttitle{Julio}{31}
        \gantttitle{Agosto}{31}
        \gantttitle{Sept.}{3} \\
        \ganttbar{F0 · Cimientos}{1}{17} \\
        \ganttbar{F1 · Embeddings y vector DB}{18}{31} \\
        \ganttbar{F2 · Pipeline RAG}{32}{48} \\
        \ganttbar{F3 · Aplicación web}{49}{59} \\
        \ganttbar{F4 · Validación y ablation}{60}{68} \\
        \ganttbar{F5 · Cierre y defensa}{69}{72} \\
        \ganttmilestone{Entrega}{72}
    \end{ganttchart}%
    }
    \caption{Planificación temporal del proyecto (del 24 de junio al 3 de septiembre de 2026).}
    \label{fig.gantt}
\end{figure}

% Sustituir esta figura por una captura real del tablero de GitHub Projects (IT-48).
